% Bibliographic records and exact source leaves checked 9 October 2026.
% Source access limits are part of the records below; this is not a global novelty audit.
\begin{thebibliography}{99}
\addcontentsline{toc}{chapter}{Bibliography}
\small
\setlength{\itemsep}{2pt}
\setlength{\parsep}{0pt}
\setlength{\parskip}{0pt}

\bibitem{Lehmer}
D. H. Lehmer, \emph{The vanishing of Ramanujan's function $\tau(n)$},
Duke Mathematical Journal \textbf{14} (1947), 429--433.
\href{https://doi.org/10.1215/S0012-7094-47-01436-1}{DOI: 10.1215/S0012-7094-47-01436-1}.
\href{https://oeis.org/A000594/a000594.pdf}{Original scan}.
The scan was authenticated, but its original proof was not readable through the source-check route. The prime reduction is a classical attributed input and is proved separately in this reader.

\bibitem{BCO}
Jennifer S. Balakrishnan, William Craig, and Ken Ono,
\emph{Variations of Lehmer's Conjecture for Ramanujan's tau-function},
\href{https://arxiv.org/pdf/2005.10345}{arXiv:2005.10345v1}, 20 May 2020.
Checked classical leaves: equation~(1.2), p.~1, for the congruences; Theorem~2.2, p.~4, for the Hecke relations and bound; the beginning of \S3, p.~7, for the Jacobi proof of the odd-square parity criterion. The omitted-value and integral-point proofs are outside this edition's source check.

\bibitem{DHZ}
Maarten Derickx, Mark van Hoeij, and Jinxiang Zeng,
\emph{Computing Galois representations and equations for modular curves $X_H(\ell)$},
\href{https://arxiv.org/pdf/1312.6819v2}{arXiv:1312.6819v2}, 18 March 2014.
Corollary~1.2, p.~2; its computational proof in \S4, p.~11.
The cited nonvanishing range has the strict endpoint
$n<816212624008487344127999$. This is a historical source result, not an edition replay or a verified latest record.

\bibitem{DeligneModular}
Pierre Deligne, \emph{Formes modulaires et repr\'esentations $\ell$-adiques},
S\'eminaire Bourbaki 1968/69, expos\'e~355, published 1971, 139--172.
\href{https://www.numdam.org/item/SB_1968-1969__11__139_0.pdf}{NUMDAM original}.
Theorem~5.6 supplies the implication from the Weil conjectures to the Ramanujan bound.
The author's \href{https://publications.ias.edu/sites/default/files/5_FormesModulaires.pdf}{English translation}, Proposition~4.8 and equation~(5.7), was also checked.
Its cohomological $F$ is inverse arithmetic Frobenius; the reader uses the rank-two arithmetic normalization specified in~\cite{SerreBourbaki}.

\bibitem{DeligneWeil}
Pierre Deligne, \emph{La conjecture de Weil: I},
Publications Math\'ematiques de l'IH\'ES \textbf{43} (1974), 273--307.
\href{https://doi.org/10.1007/BF02684373}{DOI: 10.1007/BF02684373}.
\href{https://www.numdam.org/article/PMIHES_1974__43__273_0.pdf}{Theorem~1.6, p.~276}.
For the modular specialization giving $|\tau(p)|\le2p^{11/2}$, this purity theorem is used with the implication in~\cite{DeligneModular}.

\bibitem{SerreBourbaki}
Jean-Pierre Serre, \emph{Congruences et formes modulaires (d'apr\`es H. P. F. Swinnerton-Dyer)},
S\'eminaire Bourbaki 1971/72, expos\'e~416, Lecture Notes in Mathematics \textbf{317} (1973), 319--338.
\href{https://www.numdam.org/item/SB_1971-1972__14__319_0.pdf}{NUMDAM PDF}.
\S3.1, pp.~331--332, supplies the established integral $\ell$-adic representations and their unramified Frobenius polynomials. Theorem~10 and \S3.3, p.~335, concern image exceptions. These statements do not assert joint surjectivity for an arbitrary product of reductions.

\bibitem{SerreCongruences}
Jean-Pierre Serre, \emph{Une interpr\'etation des congruences relatives \`a la fonction $\tau$ de Ramanujan},
S\'eminaire Delange--Pisot--Poitou \textbf{9} (1967/68), no.~1, expos\'e~14, 1--17.
\href{https://www.numdam.org/item/SDPP_1967-1968__9_1_A13_0.pdf}{NUMDAM PDF}.
\S2.2, equation~(21), p.~4, for modulo~691; \S2.1, equation~(16), p.~3, and \S4.5, p.~11, for the stated prime congruences modulo~23 and~$23^2$.
The representation existence is still conjectural in this 1968 text; its established version is cited from~\cite{SerreBourbaki}.

\bibitem{Han}
Guo-Niu Han, \emph{The Nekrasov--Okounkov hook length formula: refinement, elementary proof, extension and applications},
\href{https://arxiv.org/pdf/0805.1398}{arXiv:0805.1398v1}, 9 May 2008,
Definition~1.1 and Theorems~1.1--1.2, p.~3.
Published in Annales de l'Institut Fourier \textbf{60} (2010), 1--29,
\href{https://doi.org/10.5802/aif.2515}{DOI: 10.5802/aif.2515}.
The checked specialization is $t=5$, $z=25$: shell $\sum v_i^2=10n$ and Vandermonde weight $D(v)/288$. Integrality and the sign consequences are separate project arguments.

\bibitem{GKS}
Frank Garvan, Dongsu Kim, and Dennis Stanton, \emph{Cranks and $t$-cores},
Inventiones Mathematicae \textbf{101} (1990), 1--17.
\href{https://doi.org/10.1007/BF01231493}{DOI: 10.1007/BF01231493}.
Checked \href{https://www-users.cse.umn.edu/~stant001/PAPERS/cra2.pdf}{author manuscript},
\S2, Bijections~1--2; \S5, Theorem~4 and equation~(5.3), printed p.~8.
The exact five-core count is translated to V-coded packets using Han's bijection; it is not a sign assertion about the weighted sum.

\bibitem{IW}
Tomoyoshi Ibukiyama and Brandon Williams,
\emph{Hermitian modular forms and algebraic modular forms on $SO(6)$},
\href{https://arxiv.org/pdf/2505.23497}{arXiv:2505.23497v1},
header 29 May 2025, internal date 12 June 2025.
\S3.4, printed pp.~18--19, gives vector-valued harmonic theta modularity and weight $\nu+n/2$.
The scalar full-level sign-line descent is a project argument. The paper's conjectural Hermitian correspondences are not assumed.

\bibitem{Stein}
William Stein, \emph{Modular Forms: A Computational Approach},
author online chapter \href{https://wstein.org/books/modform/modform/level_one.html}{\emph{Modular Forms of Level 1}}, checked 9 October 2026.
Theorem~2.11 for valence; Theorem~2.14, Corollaries~2.15--2.16, and Theorem~2.17 for the level-one structure and dimensions.
The chapter's Eisenstein series differ by nonzero scalars from the reader's constant-one normalization; its displayed normalized identity is $\Delta=(E_4^3-E_6^2)/1728$.

\bibitem{Fonseca}
Tiago J. Fonseca,
\emph{Higher Ramanujan equations I: moduli stacks of abelian varieties and higher Ramanujan vector fields},
\href{https://arxiv.org/pdf/1612.05081}{arXiv:1612.05081v1},
header 15 December 2016, internal date 5 November 2018.
Introduction, p.~2, and Remark~6.4, p.~29, for the normalized classical differential relations. The reader's polynomial residual identity is reconstructed from these inputs, rather than inferred from bounded coefficient tests.

\bibitem{Niebur}
Douglas Niebur, \emph{A formula for Ramanujan's $\tau$-function},
Illinois Journal of Mathematics \textbf{19} (1975), 448--449.
\href{https://doi.org/10.1215/ijm/1256050746}{DOI: 10.1215/ijm/1256050746}.
The original proof could not be accessed through the check route. The universal convolution identity was checked in the primary-paper restatement~\cite{MoriyaSmyth}; the differential reconstruction in this reader uses~\cite{Fonseca,Stein}.

\bibitem{MoriyaSmyth}
B. K. Moriya and C. J. Smyth,
\emph{Index-dependent divisors of coefficients of modular forms},
\href{https://webhomes.maths.ed.ac.uk/~chris/papers/Moriya_Smyth_final.pdf}{author manuscript},
dated 1 May 2013, accepted 28 April 2013.
\S7, printed pp.~10--11, restates Niebur's identity for every $n\ge1$ and gives the symmetric half-sum convention.
\S2.1 explicitly retains the alternative $a(p)=0$ in the size/divisibility hypothesis.
The locators refer to the inspected manuscript.

\bibitem{LRTrace}
Nik Lygeros and Olivier Rozier,
\emph{A New Solution to the Equation $\tau(p)\equiv0\pmod p$},
Journal of Integer Sequences \textbf{13} (2010), article~10.7.4.
\href{https://cs.uwaterloo.ca/journals/JIS/VOL13/Lygeros/lygeros5.pdf}{Journal PDF},
\S\S4--6, printed pp.~3--6, for the trace moments, Catalan inverse, and reduced-form convention.
The Hurwitz count includes imprimitive forms and exceptional weights $1/2$ and $1/3$; $6H$ is the stored integer. The zeroth moment keeps half of the $t=0$ term.

\bibitem{LRPrime}
Nik Lygeros and Olivier Rozier,
\emph{Odd prime values of the Ramanujan tau function}, Ramanujan Journal \textbf{32} (2013), 269--280.
\href{https://doi.org/10.1007/s11139-012-9420-8}{DOI: 10.1007/s11139-012-9420-8}.
Checked \href{https://lygeros.org/wp-content/uploads/root/11713_Odd_prime_values_of_the_Ramanujan_tau_function.pdf}{author PDF},
p.~2, equations~(3)--(11), as restatements of the sharper small-prime congruences, with their stated residue and unit hypotheses. The original earlier proofs of the entire congruence table were not separately audited.

\bibitem{Chebotarev}
Peter Stevenhagen and Hendrik W. Lenstra Jr.,
\emph{Chebotar\"ev and his density theorem},
The Mathematical Intelligencer \textbf{18} (1996), no.~2, 26--37.
Checked \href{https://pub.math.leidenuniv.nl/~lenstrahw/PUBLICATIONS/1996d/art.pdf}{author scan},
p.~31 for the natural-density convention and p.~34 for the conjugacy-class theorem.
The historical project URL failed; this author-hosted copy was inspected successfully.
The theorem applies to one fixed finite quotient and a nonempty conjugacy class, giving density $|C|/|G|$.

\bibitem{CF}
Pierre Colmez and Jean-Marc Fontaine,
\emph{Construction des repr\'esentations $p$-adiques semi-stables},
Inventiones Mathematicae \textbf{140} (2000), 1--43.
\href{https://doi.org/10.1007/s002220000042}{DOI: 10.1007/s002220000042}.
Checked \href{https://webusers.imj-prg.fr/~pierre.colmez/CF.pdf}{37-page author manuscript}:
Theorem~A, p.~2; \S3.4, pp.~21--22; \S4.1 and Corollary~4.3, pp.~23--24; the period-ring identity on p.~11.
These are manuscript locators. Local realization does not identify the constructed object with the restriction of Deligne's global representation.

\bibitem{TauContext2026}
David Kurniadi Angdinata et al.,
\emph{ABC implies that Ramanujan's tau function misses almost all primes},
\href{https://arxiv.org/pdf/2603.29970}{arXiv:2603.29970v3},
25 April 2026, internal date 28 April 2026.
Only the contemporary introductory statement of Lehmer's conjecture and the classical inputs in \S2.1, equation~(2.1), and Proposition~6 were checked here. The paper's conditional prime-value theorem is not an input to this reader.

\bibitem{Withdrawn2025}
Minjia Shi, Lu Wang, and Patrick Sol\'e,
\emph{Proof of the Lehmer conjecture on Ramanujan's $\tau$ function},
\href{https://arxiv.org/abs/2503.23498}{arXiv:2503.23498v2}, 1 April 2025, withdrawn.
Only the explicit withdrawal metadata was checked. This title is not cited as an established proof or as a theorem input.

\end{thebibliography}
